\chapter{What Each Interview Tests, by Role}\label{ch:iv:what-each-interview-tests-by-role}

Two candidates leave the same building on the same afternoon with opposite verdicts on the same probability
question. The trading candidate reached the right answer in four minutes and was marked down; the developer
candidate gave a wrong first answer, tested it on a small case, found the error and fixed it, and was marked
up. Nobody was inconsistent. The two interviewers were measuring different jobs: a trader must be right fast
and say how sure she is; a developer must be right in the end and know how he knows. This chapter is about
reading which job a question is measuring, because the same question has a different best answer in each
room.

\section{What an interview measures}

An interview is a measurement. The firm wants to know how well a person will do a job it can describe; it
cannot watch the person do the job for a year, so it samples a few hours of behaviour that it believes
predicts the year. Everything in this book follows from taking that sentence seriously: a question is an
instrument, an answer is a reading, and the reading is compared with a scale that the interviewer, not the
candidate, holds.

\begin{definition}[Interview loop, technical interview]\label{def:iv:what-each-interview-tests-by-role:loop}
An \emph{interview loop}\index{interview loop} is the ordered set of interviews one candidate sits for one
role, with the rule that turns their separate scores into a decision. A \emph{technical
interview}\index{technical interview} is one whose questions have answers that can be checked: a number, a
proof, a program, a design whose properties can be argued. Its opposite is the behavioural interview of
\cref{ch:iv:behavioural-and-fit}, whose answers are accounts of past conduct.
\end{definition}

\begin{definition}[Interview rubric]\label{def:iv:what-each-interview-tests-by-role:rubric}
An \emph{interview rubric}\index{interview rubric} is the scoring sheet for one question or one interview: a
short list of dimensions (for example correctness, speed, method, communication), each with anchored levels
that say what an answer at that level contains, filled in by the interviewer before the loop's scores are
compared.
\end{definition}

The rubric is the instrument that makes an interview a structured interview in the sense of One Quant
Book~16, chapter~10: the same questions for every candidate, answers scored against anchors agreed in
advance, interviewers scoring independently before they confer. The research on selection that Book~16
summarises ranks structured interviews first among selection procedures once earlier overcorrections are
removed (Sackett, Zhang, Berry and Lievens, 2022), and the components that make an interview structured were
catalogued a quarter of a century earlier (Campion, Palmer and Campion, 1997). Two consequences matter to a
candidate. First, the rubric rewards what it lists and nothing else: an elegant digression scores nothing on
a sheet that asks for a number and a check. Second, the rubric is written for the job, so the best answer to
the same question differs by role, as the opening scene shows.

A loop is also a chain of filters, and its arithmetic is worth knowing before any question is asked.

\begin{proposition}[Pass rates through a loop]\label{prop:iv:what-each-interview-tests-by-role:funnel}
Let an application pass stage $k$ with probability $p_k$, independently of the other stages, for
$k = 1, \dots, K$. The probability that it ends in an offer is $\prod_{k=1}^K p_k$; the expected number of
stages sat per application is $\sum_{k=1}^{K} \prod_{j<k} p_j$; and with $n$ independent applications of offer
probability $q$, the chance of at least one offer is $1 - (1-q)^n$.
\end{proposition}

\begin{proof}
Stage $k$ is reached only if stages $1, \dots, k-1$ are passed, so it is sat with probability
$\prod_{j<k} p_j$; summing these indicators gives the expected count. The last statement is the complement of
$n$ independent failures.
\end{proof}

\begin{example}[A three-stage funnel]\label{ex:iv:what-each-interview-tests-by-role:funnel}
An online assessment passed three times in ten, a screen passed one time in two and a final round passed two
times in five give an offer probability of $0.3 \times 0.5 \times 0.4 = 0.06$ per application. Each
application costs on average $1 + 0.3 + 0.15 = 1.45$ stages, so an offer costs about $1.45/0.06 \approx 24.2$
stages sat and $1/0.06 \approx 16.7$ applications. Twenty applications give at least one offer with
probability $1 - 0.94^{20} \approx 0.71$. The numbers are illustrative; the lesson is structural: most of the
variation between candidates is in the first stage, which is the cheapest to improve
(\cref{ch:iv:online-assessments}).
\end{example}

The independence assumption is optimistic in one direction and pessimistic in the other. A candidate who
fails one firm's assessment is more likely to fail the next firm's, because the same weakness is measured
twice; and a candidate who has sat ten final rounds is better at the eleventh. Both effects mean that the
early applications are the ones to spend preparation on.

\section{The roles and what each is probed on}

The roles themselves are described in One Quant Book~17, chapters 16 to 25. The interviewer's question is
narrower: which abilities separate a good holder of this role from a poor one in the first two years, and
which of them can be measured in an hour? The table below is this book's answer, and it is also its map.

\begin{center}
\small
\begin{tabular}{@{}lcccccc@{}}
\toprule
Question family (chapter) & Trader & Researcher & Developer & ML engineer & Bank quant & Risk \\
\midrule
Mental arithmetic (8) & $\bullet$ & $\circ$ & & & $\circ$ & $\circ$ \\
Estimation (9) & $\bullet$ & $\circ$ & $\circ$ & $\circ$ & $\circ$ & $\circ$ \\
Probability (10, 11) & $\bullet$ & $\bullet$ & $\circ$ & $\circ$ & $\bullet$ & $\bullet$ \\
Brainteasers (12) & $\bullet$ & $\circ$ & $\circ$ & & & \\
Betting and market making (13) & $\bullet$ & $\circ$ & & & & $\circ$ \\
Statistics (14) & $\circ$ & $\bullet$ & & $\bullet$ & $\circ$ & $\bullet$ \\
Linear algebra, calculus (15) & & $\bullet$ & & $\circ$ & $\bullet$ & $\circ$ \\
Stochastic calculus (16) & & $\circ$ & & & $\bullet$ & $\circ$ \\
Options, fixed income (17, 18) & $\bullet$ & $\circ$ & & & $\bullet$ & $\bullet$ \\
Machine learning (19) & & $\bullet$ & & $\bullet$ & & \\
Research case study (20) & & $\bullet$ & & $\bullet$ & & \\
Algorithms, languages (21--24) & $\circ$ & $\circ$ & $\bullet$ & $\bullet$ & $\circ$ & \\
Systems, design (25, 26) & & & $\bullet$ & $\circ$ & & \\
SQL and data (27) & & $\circ$ & $\circ$ & $\bullet$ & & $\bullet$ \\
Behavioural (28) & $\bullet$ & $\bullet$ & $\bullet$ & $\bullet$ & $\bullet$ & $\bullet$ \\
\bottomrule
\end{tabular}

\smallskip
{\small Where each role's interviews spend
their time: $\bullet$ a core family, asked in most loops; $\circ$ asked in some. A chapter number follows each
family.}
\end{center}

Read by column, the table describes six loops.

\emph{Trading roles} (the risk trader and execution trader of One Quant Book~17, chapter~16, and the market
maker's trader of Book~11). The interview measures speed with numbers, calibrated judgement under uncertainty,
and behaviour when a counterparty knows more: mental arithmetic, probability at speed, estimation with an
interval, and games in which the candidate quotes prices and is traded against. The best answer is fast,
approximately right, and says how sure it is; a slow exact answer is scored lower than a quick answer with
the right error bar.

\emph{Research roles} (the quantitative researcher of Book~17, chapter~17). The interview measures whether the
candidate can turn a vague question into a testable one and refuse to be fooled: probability and statistics in
depth, regression pathologies, overfitting, and a research case study or take-home. The best answer states
assumptions, tests them, and says what would change the conclusion.

\emph{Developer roles} (the quant developer of Book~16, chapter~21, the research engineer and low-latency
engineer of Book~17, chapters 19 and~20). The interview measures working code, reasoning about cost, and
knowledge of what the machine does: algorithms, a language in depth, concurrency and systems, and a design
problem. The best answer is correct, tested, and explains its complexity and its failure modes; the trading
candidate's four-minute answer in the hook would score badly here if it had not been checked.

\emph{Machine-learning engineers} (Book~12, chapter~28) sit between research and development: statistics and
validation with time, model design, and production code.

\emph{Bank quants} (the desk strategist, model validator and risk quant of Book~17, chapter~18). The interview
measures mathematics of pricing and its implementation: stochastic calculus, derivatives, numerical methods and
a language. A model validator is asked more about what can go wrong with a model, a desk strategist more about
what the desk needs by tomorrow morning.

\emph{Risk and control roles} (the risk quant and the control functions of Book~17, chapter~24). The interview
measures statistics of losses, derivatives at the level of their sensitivities, and judgement about what to
escalate.

\begin{remark}[One question, several rubrics]\label{rem:iv:what-each-interview-tests-by-role:rubrics}
``How many rolls of a fair die, on average, until the first six?'' has one answer, six. A trading rubric
scores the time to answer and whether the candidate volunteers the variance (thirty) when asked how long the
wait could be. A research rubric scores the argument (the geometric law, or a one-line first-step equation).
A developer rubric scores a check: a ten-line simulation, or the observation that the answer must exceed
one and be finite. The candidate who knows the room gives the answer that rubric can record.
\end{remark}

\section{Formats and the time each question gets}

Most loops are built from a small number of formats: a timed online assessment, a screen by telephone or
video, and a final round of several interviews, sometimes with a trading game, a group exercise or a
presentation (\cref{ch:iv:the-process-by-firm-type,ch:iv:the-final-round-and-trading-games}). Within one
interview the question types have typical lengths, and the length tells the candidate what is expected.

\begin{method}[Reading a question by its time budget]\label{met:iv:what-each-interview-tests-by-role:budget}
\begin{enumerate}
\item \emph{Seconds}: arithmetic, a sign, a definition. The interviewer wants the answer, not the method.
\item \emph{One to five minutes}: a probability, an estimate, a Greek, an output to predict. Say the method in
one sentence, compute, check once, stop.
\item \emph{Ten to twenty minutes}: a derivation, a coding problem, a game. Clarify, give a baseline, improve,
test; think aloud (\cref{ch:iv:phone-and-technical-screens}).
\item \emph{Thirty minutes or more}: a design, a case study, a take-home. Requirements first, numbers early,
structure before detail (\cref{ch:iv:system-design,ch:iv:research-case-studies-and-take-homes}).
\item If the time budget is unclear, ask: ``Do you want the number, or the argument?'' costs three seconds.
\end{enumerate}
\end{method}

A loop rarely fails a candidate on one bad answer. It fails candidates whose answers are consistently in the
lower anchors on the dimensions the role weighs: the trader who never gives a number, the researcher who never
checks an assumption, the developer who never tests. \Cref{iq:iv:what-each-interview-tests-by-role:5} shows
the arithmetic of the other failure: a loop that demands a pass in every interview rejects good candidates
often, which is why well-designed loops average scores rather than multiply vetoes.

\section{Preparing by role}

This book is organised by question family, not by role. A reader preparing for one role reads the core
chapters of its column in the table of the previous section first, then the process
chapters (Part~I), then \cref{ch:iv:mock-interviews}, whose six transcripts show one complete interview for
each column. Each chapter's bank is ordered from one star to three and tagged with the roles and firm types
for which the question is typical; the tags describe the kind of question, never a firm that asked it. Every
question in the book is new: none is reprinted from another book of this series or elsewhere, and a classic
puzzle appears only as a variant whose origin is cited.

\begin{example}[A preparation order for a developer with a mathematics degree]\label{ex:iv:what-each-interview-tests-by-role:order}
The degree covers probability and linear algebra in principle, not at interview speed. The order is:
\cref{ch:iv:algorithms-and-data-structures} and the language chapter of the role (C++ or Rust), then
\cref{ch:iv:concurrency-operating-systems-and-networks,ch:iv:system-design}, then
\cref{ch:iv:probability-i} for speed, then Part~I and \cref{ch:iv:behavioural-and-fit}. The mathematics is
revised last because it decays slowest.
\end{example}

\section{Worked answers}

Each chapter of this book ends its lesson with a few questions answered in full, as a strong candidate would answer them
aloud: the reasoning first, the number, then the check. They are not in the bank, and their numbers are asserted by the
chapter's tests like every other number in the book.

\begin{example}[Why a strong candidate needs several loops]\label{ex:iv:what-each-interview-tests-by-role:several}
\emph{``You are well above the bar: your score in any one interview is your ability, one standard deviation above the pass
mark, plus independent standard normal noise. A loop has five interviews and needs all five passed. How likely is an
offer, and what does it mean for how many processes you run?''}

\emph{Answer.} One interview is passed with probability $\Phi(1) \approx 0.841$; five in a row with $0.841^5 \approx 0.42$.
A candidate well above the bar fails more loops than they pass. With three independent processes the chance of at least
one offer is $1 - 0.58^3 \approx 0.81$; at one and a half standard deviations above the bar a single loop already gives
$\Phi(1.5)^5 \approx 0.71$. Two consequences: run several processes in parallel (\cref{ch:iv:the-process-by-firm-type}),
and read a rejection as a draw from a noisy test, not as a verdict
(\cref{fig:iv:what-each-interview-tests-by-role:loops}). \emph{Check:} with an ability of zero a single
interview is a coin toss and the loop passes one time in 32, which matches $0.5^5$.
\end{example}

\begin{omfigure}
\begin{tikzpicture}
\begin{axis}[width=9.5cm,height=5.2cm,xlabel={\small ability above the bar, $\theta$ (noise standard deviations)},
  ylabel={\small chance of an offer},xmin=-1,xmax=3,ymin=0,ymax=1,
  tick label style={font=\footnotesize},grid=major,grid style={gray!20},table/col sep=comma,
  legend cell align=left,legend pos=south east,legend style={font=\scriptsize}]
\addplot[omDef,thick] table[x=theta,y=k1]{figdata/interviews/01-what-each-interview-tests-by-role/loops.csv};
\addplot[omProp,thick,dashed] table[x=theta,y=k3]{figdata/interviews/01-what-each-interview-tests-by-role/loops.csv};
\addplot[omThm,thick,densely dotted] table[x=theta,y=k5]{figdata/interviews/01-what-each-interview-tests-by-role/loops.csv};
\legend{one interview,three,five}
\end{axis}
\end{tikzpicture}

\omcaption{Chance of passing a loop that needs every interview passed, by the candidate's ability above the bar, when each
interview adds independent standard normal noise. At $\theta = 1$ a single interview is passed 84\% of the time and a
five-interview loop 42\%. Data: \texttt{fig\_iv\_loops.py}.}\label{fig:iv:what-each-interview-tests-by-role:loops}
\end{omfigure}

\begin{example}[Reading an advertisement into a plan]\label{ex:iv:what-each-interview-tests-by-role:advert}
\emph{``An advertisement for a quant developer on an execution platform asks for modern C++, Linux and networking,
Python, `strong problem-solving', and says trading knowledge is a plus. You have forty hours over four weeks. Plan
them.''}

\emph{Answer.} Map each line to the interviews it predicts: modern C++ to output-prediction and design questions
(\cref{ch:iv:cpp}); Linux and networking to the systems interview
(\cref{ch:iv:concurrency-operating-systems-and-networks}); problem-solving to the coding rounds
(\cref{ch:iv:algorithms-and-data-structures}); Python to a screen or a take-home (\cref{ch:iv:python}); the
``plus'' to one conversation about the order path (\cref{ch:iv:system-design}). Weight by how many interviews each
feeds and how far you are from speed: 12 hours of timed coding, 10 of C++, 8 of systems, 4 of design, 4 of probability
at interview speed, 2 on the behavioural stories: 40. Book the first mock interview for the end of week two, so that
weeks three and four fix what it exposes.
\end{example}

\section{Question bank}

\begin{interviewq}[$\star$ \iqroles{trader, developer} \iqfirm{any}]\label{iq:iv:what-each-interview-tests-by-role:1}
You are asked: ``A fair coin is tossed until it shows heads twice in a row; how many tosses on average?'' Say
what a trading interviewer and a developer interviewer would each be scoring, and give the answer each would
want to hear.
\end{interviewq}

\begin{interviewq}[$\star$ \iqroles{researcher} \iqfirm{systematic fund}]\label{iq:iv:what-each-interview-tests-by-role:2}
Name three things a research interview checks that a trading interview usually does not, with one question
that would check each.
\end{interviewq}

\begin{interviewq}[$\star$ \iqroles{trader, researcher, developer} \iqfirm{any}]\label{iq:iv:what-each-interview-tests-by-role:3}
Your applications pass the online assessment one time in four, the screen one time in two and the final
round two times in five, independently. What is the chance that one application ends in an offer? How many
applications do you expect to make before the first offer, and what is the chance of at least one offer from
twelve? How many would you need for a 90\% chance?
\end{interviewq}

\begin{interviewq}[$\star\star$ \iqroles{developer} \iqfirm{proprietary firm}]\label{iq:iv:what-each-interview-tests-by-role:4}
A coding interview is scored on four rows, each from 1 to 4: correctness, complexity, testing and
communication. You wrote an $O(n \log n)$ solution to a problem with an $O(n)$ answer, found your own
off-by-one error with a test you wrote, and explained as you went. Where does your answer sit on each row,
and which single change would have raised the total most?
\end{interviewq}

\begin{interviewq}[$\star\star$ \iqroles{researcher, risk} \iqfirm{any}]\label{iq:iv:what-each-interview-tests-by-role:5}
A candidate's ability is $\theta$; each of four interviews scores $\theta + Z$ with independent standard normal
$Z$ and is passed when the score is positive. For $\theta = 0.5$, what is the chance of passing all four? Of
passing at least three? What ability does a candidate need to pass all four with probability one half, and
what does this say about loops that require every interviewer's approval?
\end{interviewq}

\begin{interviewq}[$\star\star$ \iqroles{bank} \iqfirm{bank}]\label{iq:iv:what-each-interview-tests-by-role:6}
The same bank interviews for a desk strategist and a model validator on the same derivatives desk. Give two
questions that would appear in one loop and not the other, and say why.
\end{interviewq}

\begin{interviewq}[$\star\star\star$ \iqroles{trader, risk} \iqfirm{market maker}]\label{iq:iv:what-each-interview-tests-by-role:7}
Abilities in a pool of candidates have standard deviation 1; each interviewer's score is ability plus
independent noise of standard deviation 1. What is the correlation between ability and the average score of
one, two and four interviewers? Two interviewers who confer before scoring end up sharing half of their noise;
what is the correlation of their average then? What rule for the loop follows?
\end{interviewq}

\begin{interviewq}[$\star\star\star$ \iqroles{mle, researcher} \iqfirm{systematic fund}]\label{iq:iv:what-each-interview-tests-by-role:8}
You are asked to design a one-hour \omterm{def:iv:what-each-interview-tests-by-role:loop}{technical interview} for a machine-learning engineer who will put research
models into production at a systematic fund. What do you test, in what order, and what does your rubric
record?
\end{interviewq}

\begin{omsources}
\item P.~R.~Sackett, C.~Zhang, C.~M.~Berry and F.~Lievens, ``Revisiting meta-analytic estimates of validity in
personnel selection'', \emph{Journal of Applied Psychology} 107(11), 2022, 2040--2068.
\item F.~L.~Schmidt and J.~E.~Hunter, ``The validity and utility of selection methods in personnel psychology'',
\emph{Psychological Bulletin} 124(2), 1998, 262--274.
\item M.~A.~Campion, D.~K.~Palmer and J.~E.~Campion, ``A review of structure in the selection interview'',
\emph{Personnel Psychology} 50(3), 1997, 655--702.
\item One Quant Book~16, chapter~10 (structured interviews and the hiring funnel); One Quant Book~17,
chapters 16--25 (the roles).
\end{omsources}
